\section{ねじれを保つ compact nodal block の構成}
\label{lc:construction}

本節では、固定曲線の小さい法線近傍に、\emph{compact な fiber 全体}を持つ
Lagrangian fibration を構成する。四次元の基本模型そのものは
Auroux \cite[\S5.1, Proposition 5.2]{auroux-local} の既知の模型である。
ここで必要なのは、それを曲線上の line bundles のねじれとともに扱うこと、
compact な saturated block を実際の法線近傍内に収めること、
および相対版の入力条件を正確に区別することである。
新規性の主張は行わない。

\subsection{Dual connection による coupling}

$(\bar C,\omega_C)$ を compact な symplectic surface with boundary とし、
$E_+,E_-$ を Hermitian complex line bundles とする。
以下では unitary trivialization
\begin{equation}\label{lc:construction-tensor}
 s:E_+\otimes E_-\xrightarrow{\sim}\bar C\times\CC
\end{equation}
を\emph{一つ固定する}。したがって $E_-\simeq E_+^*$ であるが、
$E_+$ 単独の自明化は指定しない。
特に、境界 frame を固定した $E_+$ の相対 Chern 数の消失は仮定しない。

$E_+$ の unitary connection を取り、$E_-$ には
式\eqref{lc:construction-tensor}に関して dual な接続を入れる。
その angular connection forms を $\alpha_+,\alpha_-$ とする。
角座標は $\theta_\pm\in\RR/\ZZ$ とし、周期は1に統一する。
局所 unitary frame で
\[
 \alpha_+=d\theta_++A,\qquad
 \alpha_-=d\theta_--A
\]
と書ける。曲率を $\kappa=dA$ とする。
射影を $p_C:E_+\oplus E_-\to\bar C$ とし、fiber coordinates $(y,z)$ に対して
\begin{equation}\label{lc:construction-coordinates}
 \rho_+=\pi|y|^2,\qquad \rho_-=\pi|z|^2,\qquad
 J=\rho_+-\rho_-,\qquad w=s(y\otimes z)
\end{equation}
と置く。この式の $\pi$ は円周率である。
次の2形式を考える。
\begin{equation}\label{lc:construction-coupling}
 \Omega_0=p_C^*\omega_C+
 d(\rho_+\alpha_++\rho_-\alpha_-).
\end{equation}

\begin{lemma}\label{lc:construction-coupling-lemma}
式\eqref{lc:construction-coupling}は zero section を越して滑らかで、
十分小さい法線近傍上で symplectic である。
接続による水平・垂直分解に関して mixed term はなく、垂直部分は
標準形 $\omega_{\CC^2}$、水平部分は $\omega_C+J\kappa$ である。
周期1の作用
\[
 (y,z)\longmapsto(e^{2\pi it}y,e^{-2\pi it}z)
\]
の Hamiltonian は $J$ であり、規約は $\iota_{X_J}\Omega_0=-dJ$ である。
\end{lemma}

\begin{proof}
一つの複素座標 $y=u+iv$ について
$\rho\,d\theta=(u\,dv-v\,du)/2$ なので、この1形式は原点を越して滑らかである。
また $\rho A$ も滑らかである。したがって式の angular notation による
見かけの特異性は実際には存在しない。
局所表示は
\[
 \Omega_0=\omega_C+\omega_{\CC^2}+d(JA).
\]
水平 lift は $\alpha_+=\alpha_-=0$ により定まり、
$\rho_\pm$ は水平輸送で一定である。
この分解で垂直部分は $\omega_{\CC^2}$、水平部分は
$\omega_C+JdA$、mixed term は零となる。
$\bar C$ は compact なので、ある $m_0>0$ に対し
$|J|<m_0$ なら $\omega_C+J\kappa$ は正の面積形式である。
これにより非退化性が従う。
最後に作用の生成ベクトル場では
$\alpha_+(X_J)=1$, $\alpha_-(X_J)=-1$ であり、
$\rho_\pm$ は不変だから、直接の contraction は $-dJ$ となる。
\end{proof}

\subsection{Compact な四次元 nodal fiber}

$c>0$ を選び、
\begin{equation}\label{lc:construction-q}
 Q=|w-c|^2
\end{equation}
と置く。$c$ は tensor trivialization に関する複素平面上の正の実数である。
$h:\bar C\to\RR$ を、境界の各 collar で正則な Morse function とする。

\begin{proposition}\label{lc:construction-commuting}
$\Omega_0$ に関して $J,Q,h\circ p_C$ は Poisson commute する。
この結論は $\kappa=0$ を必要としない。
\end{proposition}

\begin{proof}
$J$ と $w$ は水平輸送で一定なので、それらの微分の水平部分は零である。
したがって $X_J,X_Q$ は垂直である。
一方 $h\circ p_C$ の微分は垂直部分を持たないので、その Hamiltonian vector field
は水平である。水平・垂直部分の symplectic orthogonality から
$\{J,h\}=\{Q,h\}=0$ を得る。
$w=yz$ は逆位相の circle action に不変なので $\{J,Q\}=0$ でもある。
\end{proof}

$U$ を $(J,Q)=(0,c^2)$ を内部に含む小さい closed disk とし、
ある $m>0$ と $0<\eta<c^2$ に対して
\begin{equation}\label{lc:construction-window}
 U\subset\{\,|J|\le m,\ |Q-c^2|\le\eta\,\}
\end{equation}
と仮定する。必要なら長方形の角を丸めて $U$ を選んでもよい。
$M_U=(J,Q)^{-1}(U)$ と置く。
ここでは全ての $p\in\bar C$ 上の fiber を含める。

\begin{lemma}\label{lc:construction-radius}
$M_U$ は compact であり、その全ての点は
\begin{equation}\label{lc:construction-radius-formula}
 |y|^2+|z|^2
 =\sqrt{(J/\pi)^2+4|w|^2}
 \le\sqrt{(m/\pi)^2+4\bigl(c+\sqrt{c^2+\eta}\bigr)^2}
\end{equation}
を満たす。
$m=O(c)$, $\eta=\vartheta c^2$ $(0<\vartheta<1)$ とすれば、
$c\to0$ とともに $M_U$ 全体が任意に小さい一様な法線 disk bundle 内に入る。
\end{lemma}

\begin{proof}
$|w-c|=\sqrt Q$ なので $|w|\le c+\sqrt Q$ である。
$a=|y|^2$, $b=|z|^2$ と置くと、
$a-b=J/\pi$, $ab=|w|^2$ だから
$(a+b)^2=(J/\pi)^2+4|w|^2$ を得る。
これが所要の評価である。
$\bar C$ と $U$ は compact で、上の評価は全 fiber を一様に有界な
disk bundle 内に置く。$M_U$ はその中の閉集合なので compact である。
最後の主張も右辺から直ちに従う。
\end{proof}

\begin{lemma}\label{lc:construction-normal-fibers}
各 $p\in\bar C$ の上で、$U$ に制限した $(J,Q)$ の正則 fiber は $T^2$、
$(J,Q)=(0,c^2)$ の fiber は一度 pinch した $T^2$ である。
この特異 fiber はただ一つの focus--focus critical point を持つ。
\end{lemma}

\begin{proof}
$\CC^2$ の逆位相 circle action による quotient は
$(J,w)\in\RR\times\CC$ で記述できる。
実際、$J,w$ から $|y|,|z|$ は一意に決まり、$w\ne0$ なら
二つの位相の和が決まり、残る位相差が一つの circle orbit を与える。
$w=0,J\ne0$ でも一方の座標は非零であり、orbit は circle である。
ただし $J=w=0$ の orbit は原点一つである。

$(J,Q)$ を固定すると $w$ は $|w-c|=\sqrt Q$ という円を走る。
この円が collapse point $(J,w)=(0,0)$ を通らなければ、
fiber は円上の principal circle bundle である。
そのような束は自明なので fiber は $T^2$ となる。
$J=0,Q=c^2$ では円が $w=0$ を一度通り、ちょうど一つの orbit circle が潰れる。
これは一度 pinch した torus である。

臨界点も直接確認できる。
原点以外で $dQ=\lambda\,dJ$ とすると、複素微分の二式にそれぞれ
$y,z$ を掛けることにより $\lambda(|y|^2+|z|^2)=0$ が従う。
したがって $\lambda=0$ であり、$w=c$、すなわち $Q=0$ が必要となる。
この locus は $\eta<c^2$ によって $U$ の外にある。
原点では
\[
 Q-c^2=-2c\operatorname{Re}(yz)+|yz|^2.
\]
従って二次部分は
$\bigl(\pi(|y|^2-|z|^2),-2c\operatorname{Re}(yz)\bigr)$ であり、
標準的な非退化 focus--focus pair と線形 symplectic 同値である。
この式は Auroux の模型の局所式でもある\cite{auroux-local}。
\end{proof}

\subsection{Pants と全 punctured curve 上の fibration}

$G_h$ を $h$ の connected level components による Reeb graph、
$r_h:\bar C\to G_h$ を quotient map とする。
Critical values は互いに異なるとし、各 critical level component が
一つの figure-eight である Morse saddle のみを考える。

\begin{theorem}\label{lc:construction-block}
$U$ を十分小さく選ぶと
\[
 F=(J,Q,r_h\circ p_C):M_U\longrightarrow U\times G_h
\]
は proper で connected fibers を持つ stratified Lagrangian fibration である。
正則 fiber は $T^3$ であり、特異 fiber の位相型は次の通りである。
\begin{center}
\begin{tabular}{ll}
normal node、$h$ 正則 & $(\text{pinched }T^2)\times S^1$ \\
normal 正則、$h$ saddle & $T^2\times\figeight$ \\
normal node、$h$ saddle & $(\text{pinched }T^2)\times\figeight$.
\end{tabular}
\end{center}
最後の型の Euler 数は $-1$ である。
一つの pants では基底は $U\times Y$ となる。
\end{theorem}

\begin{proof}
Compactness により properness は補題\ref{lc:construction-radius}から従う。
$G_h$ の一点の逆像は、正則値なら円、頂点なら figure-eight である。
どちらの上でも $E_+$ の circle bundle は位相的に自明である。
したがって補題\ref{lc:construction-normal-fibers}の normal fiber との積として
全 fiber の位相型が得られ、特に連結である。
これは位相的な積の主張であり、接続が symplectic な積であるとは主張しない。

正則 locus では三つの可換な積分が独立であり、fiber の次元は3なので
Lagrangian である。特異 fiber については、connection splitting を使う。
水平射影は $h$ の level graph の滑らかな stratum に接し、次元は高々1である。
従って水平の2形式は零になる。
垂直成分は四次元の regular torus または pinched torus の isotropic stratum に接する。
Mixed term も零だから、全ての fiber strata は isotropic となる。
各 fiber には稠密な3次元 strata がある。

Saddle と node が同時に起こる点では、zero section 上の symplectic form は
$\omega_C\oplus\omega_{\CC^2}$ であり、二次部分は Morse saddle と
focus--focus pair の直和になる。従って Williamson 型は
hyperbolic block 一つと focus--focus block 一つである。
位相的積から Euler 数は $1\cdot(-1)=-1$ である。
\end{proof}

\begin{remark}\label{lc:construction-branched}
この基底は普通の3-ball 内に thin graph discriminant を持つ模型ではない。
$U\times Y$ には分岐 sheet があり、外周 $\partial U$ と pants の三つの端もある。
これらの外側境界を消す大域接着は本定理の結論に含まれない。
\end{remark}

\begin{corollary}\label{lc:construction-free-curve}
$\bar C$ の genus が $g$、境界成分数が $n\ge2$ であるとする。
端の fibration を自由に選べるなら、tensor bundle が自明なとき、
$2g+n-2$ 個の Euler 数 $-1$ の fiber を持つ上記の局所 block を構成できる。
境界のある曲面では $H^2(\bar C;\ZZ)=0$ なので、tensor bundle の絶対的な自明性は
自動的に成り立つ。
\end{corollary}

\begin{proof}
一つの境界成分を incoming、残り $n-1$ 個を outgoing とする。
$\bar C$ はこの cobordism として、$2g+n-2$ 個の index 1 handles のみで構成できる。
例えば genus を作る各 handle pair に二つの saddles を用い、
残り $n-2$ 個の saddles で outgoing circles を増やす。
この handle decomposition から、内部の極値を持たず、critical values が相異なる
Morse function $h$ を得る。各 critical level component は figure-eight である。
Tensor trivialization と任意の unitary connection を選び、定理を適用する。
相対的に指定された境界 frame の延長は、この証明では要求していない。
\end{proof}

\subsection{Actual symplectic neighborhood への移送と相対版}

ここでは germs の同型と、compact な nodal fiber を含む saturated neighborhood の同型を区別する。
Focus--focus point の局所標準形だけでは、後者の同一視は得られない。
実際、semiglobal な singular foliation には return-action の情報がある
\cite[Theorem 2.1]{vungoc-local}。

\begin{proposition}\label{lc:construction-relative-moser}
$(M,\Omega_{\rm act})$ の中に $(\bar C,\omega_C)$ が symplectic に埋め込まれ、
その symplectic normal bundle が $E_+\oplus E_-$ と同定されているとする。
境界の外側へ小さい collar と幾何データを延長してから近傍定理を適用するものとする。
境界の互いに素な collar 上で、式\eqref{lc:construction-coupling}による
\emph{実際の} symplectic tubular identifications が与えられているとする。
次の compatibility を仮定する。
\begin{enumerate}
\item End tensor trivializations は一つの $s$ に延びる。
これは tensor line の、指定 end frames に関する相対 $c_1$ の消失と同値である。
\item End connections はその $s$ に関して dual であり、与えた normal bundle の接続として扱う。
\item End Morse germs は一つの所定の $h$ に延びる。
End nodal maps は、同じ十分小さい $c$ による $(J,Q)$ そのものである。
\item End tubular maps の first jets は上記の一つの symplectic normal identification と一致する。
\end{enumerate}
このとき小さい一様な法線近傍上に、指定 collar で与えた写像と一致する
symplectic embedding が存在する。
さらに $c,m,\eta$ を十分小さく選べば、\emph{全ての compact fibers を含む}
$M_U$ をその embedding で $M$ に移せる。
既存の circle action があり、入力同定が equivariant なら、移送も equivariant にできる。
\end{proposition}

\begin{proof}
接続全体の空間は affine なので、$E_+$ の end connections を
partition of unity により大域接続へ延ばし、$E_-$ には dual connection を入れる。
これにより小さい指定 collar 上の接続そのものを保つ。
相対 tubular neighborhood construction により、所定 end maps と一致し、
zero section 上の first jet が指定された normal identification である滑らかな
embedding $\Phi_0$ を取る。
この相対構成は次のようにも確認できる。
最初に任意の tubular map を取り、end map との差を考える。
差は zero section とその first jet を固定する germ である。
法線方向の dilation で conjugate することで、この差を恒等写像へ isotope できる。
その isotopy の vector field を、互いに素な end collars の少し大きい近傍で cutoff し、
積分して最初の tubular map を補正すればよい。
Compact group の作用がある場合は、この全構成を invariant に選ぶ。

$\Omega_1=\Phi_0^*\Omega_{\rm act}$ とし、$\Delta=\Omega_1-\Omega_0$ と置く。
二つの形式は zero section 上で\emph{全ての tangent vectors に対して}一致し、
指定 end tubes では近傍全体で一致する。
法線方向の radial contraction の homotopy operator によって
\[
 \Delta=d\zeta,\qquad \zeta|_{\bar C}=0,
 \qquad \zeta=0\quad\text{on the smaller end tubes}
\]
を得る。Contraction は base point を動かさないので、最後の相対条件も保たれる。
必要なら compact group で平均して $\zeta$ を invariant にする。

$\Omega_t=\Omega_0+t\Delta$ は、compact な $\bar C$ の十分小さい一様近傍上で
全ての $t\in[0,1]$ に対して symplectic である。
Moser equation $\iota_{V_t}\Omega_t=-\zeta$ を解くと、$V_t$ は zero section と
指定 end tubes で零となる。
近傍をさらに縮めれば、その flow は時刻1まで存在する。
得られた flow を $\Phi_0$ と合成すれば、所要の相対 symplectic embedding となる。
これは relative Moser の標準的証明\cite[Theorem 7.4]{cds-local}に
collar 条件を保たせたものであり、絶対 cohomology class の一致だけを使った議論ではない。

最後に補題\ref{lc:construction-radius}を用い、$M_U$ 全体がこの一様近傍に入るように
$c,m,\eta$ を選ぶ。
これにより局所二次模型の非コンパクトな部分だけでなく、閉じた nodal torus 全体が移る。
Equivariance がある場合、二つの moment maps の差は定数であり、
固定曲線上の規格化によってその定数も決まる。
\end{proof}

\begin{remark}\label{lc:construction-relative-scope}
個別の $E_+,E_-$ の相対 Chern 数が逆符号で非零でも、tensor frame が延びれば
本命題は適用できる。個々のねじれは dual connection の中に残している。
一方、任意に先に選ばれた nodal end foliations が本命題の入力になるとは主張しない。
その完全な semiglobal data が上記 coupling model と一致する同定、または別の相対補間が必要である。
また、$c$ が先に固定されている場合は、その compact nodal block 全体が
相対 Moser の定義域に入ることを別に確認する必要がある。
既存の torus を「germ だから小さくできる」として縮小することは許されない。
\end{remark}

\begin{remark}\label{lc:construction-variable-c}
$c$ を、$r_h$ に関して basic な正の滑らかな関数とし、各 saddle の近傍で定数とする
拡張も可能である。$dc$ は $h$ の Hamiltonian vector field を消すため、
$\{Q,h\}=0$ は保たれる。各 connected level 上で $c$ は一定なので fiber の位相も同じである。
これはこの特定の模型の scale を端ごとに調整する方法であり、
任意の semiglobal focus--focus invariant を補間する定理ではない。
\end{remark}
